% ==============================================================================
% PHẦN 2: TÓM TẮT BÁO CÁO (ABSTRACT)
% ==============================================================================

\chapter*{TÓM TẮT}
\addcontentsline{toc}{chapter}{Tóm tắt}

Trong bối cảnh nghiên cứu khoa học thuộc lĩnh vực Trí tuệ Nhân tạo và Khoa học Dữ liệu (AI/DS) bùng nổ theo hàm mũ với hàng nghìn ấn phẩm công bố mỗi tháng trên các kho tiền ấn phẩm (arXiv, OpenReview, CVF), việc theo dõi xu hướng, tổng hợp tài liệu chuyên sâu và truy xuất tri thức chính xác trở thành một thách thức lớn. Các mô hình ngôn ngữ lớn (LLMs) thương mại hiện nay thường xuyên gặp phải hiện tượng \textit{ảo giác học thuật} (hallucination), trích dẫn sai lệch nguồn gốc, hoặc bỏ sót cấu trúc công thức toán học do văn bản PDF học thuật bị phân mảnh. 

Đồ án này nghiên cứu và triển khai hệ thống \textbf{Khai thác Dữ liệu Nghiên cứu Khoa học Thời gian thực hướng tới Xây dựng Hệ thống Truy xuất Tri thức Tăng cường (Advanced RAG) cho miền AI/DS}. Hệ thống được thiết kế trên nền tảng kiến trúc \textbf{Medallion Lakehouse} phân tán (Bronze $\to$ Silver $\to$ Gold) kết hợp bộ lưu trữ cột \textbf{DuckDB OLAP}, hệ quản trị vector \textbf{Qdrant (HNSW)}, hệ thống điều phối \textbf{Apache Airflow 2.10}, và hàng đợi thông điệp chịu lỗi \textbf{RabbitMQ}.

Trọng tâm của công trình là nghiên cứu và thực thi \textbf{4 Trụ cột Khai phá Dữ liệu (4 Data Mining Pillars)} tích hợp trực tiếp làm bộ não điều hướng suy luận cho hệ thống RAG:
\begin{enumerate}[label=(\arabic*), itemsep=1pt, topsep=1pt, parsep=0pt]
  \item \textbf{Khai phá Tập phổ biến \& Luật kết hợp (FP-Growth)}: Khai phá 22 luật kết hợp liên ngành không tầm thường từ hơn 11,400 giao dịch (tiêu biểu như luật $\{\text{cs.AI}, \text{cs.CV}\} \Rightarrow \{\text{cs.LG}\}$ với $\text{Lift} > 2.5$), phục vụ cơ chế tự động mở rộng ngữ cảnh truy vấn (\textit{rule-based query expansion}).
  \item \textbf{Phân cụm Ngữ nghĩa Không giám sát (Clustering)}: Phân tích không gian vector 768 chiều (L2-normalized) kết hợp Centered PCA và MiniBatchKMeans ($K=6$), đối chiếu với thuật toán mật độ DBSCAN, làm sáng tỏ tính liên ngành và ranh giới mở của tài liệu AI hiện đại.
  \item \textbf{Khai phá Mạng lưới Khoa học \& Đồ thị Trích dẫn}: Dựng đồ thị có hướng 8,892 bài báo và 86,295 liên kết trích dẫn nội bộ, tính toán độ uy tín học thuật bằng thuật toán \textbf{Directed PageRank} ($\alpha=0.85$) và phân rã 92 cộng đồng trường phái nghiên cứu bằng thuật toán Louvain. Điểm PageRank được chuẩn hóa thành bộ nhân uy tín (\textit{Authority Boost}), tăng tới $+25\%$ điểm xếp hạng cho các công trình nền tảng kinh điển.
  \item \textbf{Động lượng Xu hướng \& Phát hiện Dị biệt (Trend \& Anomalies)}: Đo lường vận tốc xu hướng chuẩn hóa theo tỷ trọng thị trường ($\Delta \text{Share}$ Momentum) và áp dụng \textbf{Isolation Forest} với chuẩn hóa $\log(1+x)$ phát hiện 260 bài báo dị biệt cấu trúc (chuyên khảo trên 70,000 từ, công trình chứa hơn 3,700 công thức toán).
\end{enumerate}

Đánh giá thực nghiệm bằng bộ công cụ học thuật \textbf{DeepEval 2.x} trên tập câu hỏi vàng khẳng định hệ thống RAG cải tiến (kết hợp Cross-Encoder reranker \texttt{ms-marco-MiniLM-L-6-v2}, Reciprocal Rank Fusion, PageRank boost và negative guardrails) vượt trội rõ rệt so với Dense Retrieval đơn thuần: \textbf{Contextual Precision tăng vọt từ $0.637$ lên $0.860$ (+35.0\%)}, \textbf{Contextual Recall đạt $1.000$ (100\%)}, và \textbf{Answer Relevancy tăng $+9.4\%$ lên $0.752$}. Thử nghiệm trên phần cứng máy trạm GPU NVIDIA RTX 4060 với CUDA 12.4 ghi nhận tốc độ sinh token đạt $46.33\text{ tokens/s}$, thời gian phản hồi token đầu tiên (TTFT) chỉ $609.8\text{ ms}$ (nhanh hơn 30.7 lần so với CPU/Vulkan baseline). 

Công trình chứng minh khả năng kết hợp hài hòa giữa kỹ thuật Khai phá Dữ liệu truyền thống và Trí tuệ Nhân tạo tạo sinh hiện đại, tạo ra một giải pháp nghiên cứu khoa học có độ chính xác cao và khả năng ứng dụng thực tiễn mạnh mẽ.

\vspace{0.4cm}
\noindent\textbf{Từ khóa (Keywords):} Khai phá Dữ liệu, Nghiên cứu Khoa học, Medallion Lakehouse, Truy xuất Tri thức Tăng cường (RAG), FP-Growth, PageRank, Isolation Forest, Qdrant, AI/DS.

\clearpage

% --- ABSTRACT (ENGLISH) ---
\chapter*{ABSTRACT}
\addcontentsline{toc}{chapter}{Abstract (English)}

Amid the exponential growth of scientific literature in Artificial Intelligence and Data Science (AI/DS), researchers face unprecedented difficulties in tracking emergent trends, extracting mathematical formulas accurately, and retrieving authoritative knowledge. Commercial Large Language Models (LLMs) often suffer from severe academic hallucinations, unverified citations, and loss of LaTeX mathematical formatting due to fragmented PDF extraction pipelines.

This project designs and implements an end-to-end \textbf{Real-Time Scientific Data Mining and Advanced Retrieval-Augmented Generation (RAG) System for AI/DS}. Built upon a distributed \textbf{Medallion Lakehouse architecture} (Bronze $\to$ Silver $\to$ Gold), the system synergistically combines columnar \textbf{DuckDB OLAP}, high-dimensional \textbf{Qdrant Vector Database (HNSW)}, \textbf{Apache Airflow 2.10} orchestrator, and a fault-tolerant \textbf{RabbitMQ} message broker.

The core contribution lies in operationalizing \textbf{4 Data Mining Pillars} as the reasoning and navigational backbone for neural retrieval: (1) \textbf{Association Rule Mining (FP-Growth)} discovering non-trivial interdisciplinary patterns across 11,400+ transactions for query expansion; (2) \textbf{Unsupervised Semantic Clustering (K-Means \& DBSCAN)} revealing 6 cohesive research frontiers and open boundary distributions; (3) \textbf{Directed Citation Graph Mining} employing Directed PageRank ($\alpha=0.85$) and Louvain community detection across 86,295 internal citations to inject an academic authority multiplier (+25\%) into retrieval ranking; and (4) \textbf{Share-Normalized Trend Velocity and Isolation Forest} isolating structural outliers (theoretical monographs exceeding 3,700 formulas and 74,000 words).

Extensive empirical evaluation via \textbf{DeepEval 2.x} demonstrates that our upgraded multi-stage RAG architecture achieves an outstanding \textbf{+35.0\% improvement in Contextual Precision (from 0.637 to 0.860)}, \textbf{100\% Contextual Recall (1.000)}, and superior citation grounding. On local hardware (NVIDIA RTX 4060 GPU with CUDA 12.4), token generation throughput reaches $46.33\text{ tok/s}$ with a Time To First Token (TTFT) of $609.8\text{ ms}$ (a 30.7x speedup over CPU/Vulkan baseline). The system presents a validated blueprint for uniting classical Data Mining algorithms with modern generative AI.

\vspace{0.4cm}
\noindent\textbf{Keywords:} Data Mining, Scientific Literature, Medallion Lakehouse, Retrieval-Augmented Generation (RAG), FP-Growth, PageRank, Isolation Forest, Qdrant Vector DB, AI/DS.
